\documentclass[twocolumn]{iopjournal}

\usepackage{amsmath,amssymb,amsfonts}

\fancyhfoffset[L]{0mm}
\setlength{\marginparwidth}{0pt}
\setlength{\marginparsep}{0pt}

% ==========================================
% ZENODO PREPRINT HEADERS (ACTIVE)
% ==========================================
\renewcommand{\articletype}[1]{%
  {\vspace*{-8mm}\noindent \Large \sf Preprint}\par
  \vspace*{6mm}\noindent{\scriptsize \sf{\bfseries \MakeUppercase{#1}}}}

\fancyhead[L]{}
\fancyhead[R]{}

\begin{document}

\articletype{Supplementary Material}

\title{Supplementary Material: The Constitutive Law\\from Equilibrium Field Closure}

\author{Stephen Nelson}

\affil{Field Dynamics Research}

\ead{stephen@fielddynamics.org}

\begin{abstract}
The AQUAL field equation
$\nabla \cdot [\mu(|\nabla\Phi|/a_0)\,\nabla\Phi] = 4\pi G\rho$
requires an interpolating function $\mu(x)$ that has remained
undetermined for forty years. We derive it from three physical
constraints within the minimal single-channel saturating closure
(MSC) class. At each point, the gravitational field partitions
into a resolved fraction tracking Newtonian demand and an
unresolved fraction at saturation. Detailed balance between
resolution and relaxation propensities yields a unique equilibrium
occupancy. The deep-field scaling constraint forces the driving
exponent to unity. The result is $\mu(x) = x/(x+1)$. This
derivation uses no coupling polynomial, no M\"{o}bius map, and
no tensor rank counting. It requires only the AQUAL field equation
and three empirically grounded postulates about how a
capacity-limited field must respond to gravitational demand.
\end{abstract}

% ----------------------------------------------------------------------

\section*{The Open Problem}

Bekenstein and Milgrom~\cite{BekensteinMilgrom1984} showed that the
AQUAL field equation
\begin{equation}
  \nabla \cdot \!\left[\mu\!\left(\frac{|\nabla\Phi|}{a_0}\right)
  \nabla\Phi\right] = 4\pi G\rho
  \label{eq:aqual}
\end{equation}
derives from an action principle and therefore conserves energy,
momentum, and angular momentum. The interpolating function $\mu(x)$,
where $x = |\nabla\Phi|/a_0$, must satisfy two asymptotic limits:
\begin{align}
  \mu(x) &\to 1 \quad \text{as } x \to \infty
    \quad \text{(Newtonian recovery)}, \label{eq:P1} \\
  \mu(x)/x &\to 1 \quad \text{as } x \to 0
    \quad \text{(deep-field scaling)}. \label{eq:P2}
\end{align}
Infinitely many functions satisfy both
limits~\cite{FamaeyBinney2005,Milgrom1983}. The ``simple'' form
$x/(x+1)$, the ``standard'' form $x/\sqrt{1+x^2}$, and
transcendental forms such as $1 - e^{-\sqrt{x}}$ all produce flat
rotation curves. Milgrom's original framework does not select among
them. This note closes the gap.

% ----------------------------------------------------------------------

\section*{Three Postulates}

\paragraph{Postulate~1 (Newtonian correspondence):}
At high accelerations ($x \gg 1$), the gravitational field is fully
resolved and the constitutive response approaches unity:
$\mu(x) \to 1$.

\paragraph{Postulate~2 (finite encoding capacity):}
The gravitational field possesses finite local capacity to encode
gradients. Below a characteristic scale $a_0$, gradients become
indistinguishable from background. In this regime ($x \ll 1$), the
response becomes scale-invariant: $\mu(x)/x \to 1$.

\paragraph{Postulate~3 (single-channel saturation):}
The field operates as a single-channel processor with finite
capacity. At each point, the response partitions into a resolved
fraction $\epsilon$ (tracking Newtonian demand) and an unresolved
fraction $1 - \epsilon$ (saturated). The equilibrium partition is
determined by detailed balance.

% ----------------------------------------------------------------------

\section*{The Minimal Closure Class}

The minimal single-channel saturating closure (MSC) is defined by
five axioms, each chosen as the simplest option consistent with
the postulates:

\begin{enumerate}
\item Single occupancy variable $\epsilon \in [0,1]$.
\item Detailed balance: resolution and relaxation propensities
  are equal at equilibrium.
\item Linear saturation: the capacity factor is $(1 - \epsilon)$.
\item Monomial driving: the driving intensity is $x^n$ with
  $n > 0$.
\item Constitutive identification: $\mu(x) \equiv \epsilon(x)$.
\end{enumerate}

Each axiom is the leading-order effective field theory choice.
Multi-channel closures, nonlinear saturation, and non-monomial
driving all introduce additional parameters that current
observations do not require. Departures from MSC are testable
through transition-shape residuals in rotation curve
fits~\cite{FamaeyBinney2005}.

% ----------------------------------------------------------------------

\section*{Resolution and Relaxation}

At each point, two opposing propensities govern the field's
response:

\paragraph{Resolution propensity:}
\begin{equation}
  W_+ = \kappa_+ \cdot x^n \cdot (1 - \epsilon).
  \label{eq:Wplus}
\end{equation}
The factor $x^n$ drives resolution in proportion to the
gravitational intensity. The factor $(1 - \epsilon)$ enforces
capacity: as the channel fills, the ability to resolve more
diminishes.

\paragraph{Relaxation propensity:}
\begin{equation}
  W_- = \kappa_- \cdot \epsilon.
  \label{eq:Wminus}
\end{equation}
First-order relaxation: resolved states equilibrate with no
persistence mechanism.

% ----------------------------------------------------------------------

\section*{Deriving the Exponent}

Before solving for $\epsilon$, Postulate~2 determines the
driving exponent $n$.

At equilibrium ($W_+ = W_-$), solving for $\epsilon$:
\begin{equation}
  \epsilon = \frac{K x^n}{1 + K x^n},
  \label{eq:epsilon_general}
\end{equation}
where $K = \kappa_+/\kappa_-$. In the small-$x$ limit:
$\epsilon \approx K x^n$, so $\mu \approx K x^n$ and
$\mu/x \approx K x^{n-1}$.

Postulate~2 requires $\mu/x \to \text{const}$ as $x \to 0$.
This forces
\begin{equation}
  n = 1.
  \label{eq:n_equals_1}
\end{equation}
Any $n > 1$ gives $\mu/x \to 0$ (violating deep-field scaling).
Any $n < 1$ gives $\mu/x \to \infty$ (unbounded response). The
exponent is uniquely selected.

\paragraph{Corroborating check: flat rotation curves.}
In spherical symmetry with $\mu(x) \sim x^n$ at small $x$:
$x^n \cdot g = GM/r^2$, so $g^{n+1} = a_0^n \cdot GM/r^2$.
For circular orbits $v^2 = gr$:
$v^{2(n+1)} = a_0^n \cdot GM \cdot r^{n-1}$.
Flat rotation curves ($v$ independent of $r$) require $n = 1$.
The same exponent that satisfies Postulate~2 also produces the
observed baryonic Tully-Fisher relation $v^4 = GMa_0$.

% ----------------------------------------------------------------------

\section*{The Unique Equilibrium}

With $n = 1$ established, the equilibrium occupancy is:
\begin{equation}
  \epsilon = \frac{K x}{1 + K x}.
  \label{eq:epsilon}
\end{equation}

\paragraph{Scale absorption.}
The theory contains a single empirical scale $a_0$. The constant
$K = \kappa_+/\kappa_-$ appears only in the combination $Kx$
and can be absorbed by redefining $a_0 \to a_0/K$. This shifts
the fitted value of $a_0$ but introduces no new parameter.
Setting $K = 1$ in rescaled units:
\begin{equation}
  \boxed{\mu(x) = \frac{x}{x + 1}.}
  \label{eq:constitutive}
\end{equation}

\paragraph{Verification of Postulate~1.}
As $x \to \infty$: $\mu(x) = x/(x+1) \to 1$. Newtonian gravity
is recovered exactly.

\paragraph{Stability.}
The equilibrium is globally attracting. Define the net rate
$\dot{\epsilon} = W_+ - W_- = \kappa_+(x(1-\epsilon) -
\epsilon/K)$. This is affine in $\epsilon$ with strictly negative
slope $-\kappa_+(x + 1/K)$. The equilibrium is the unique zero
and all trajectories converge to it.

% ----------------------------------------------------------------------

\section*{Relation to the Other Routes}

The main paper derives $\mu(x) = x/(x+1)$ by two additional
routes:

\vspace{6pt}
\noindent
\begin{tabular*}{\columnwidth}{@{\extracolsep{\fill}}lp{0.37\columnwidth}@{}}
\hline
Route & Method \\
\hline
1 (companion note) & Rank-level scan of the tetrad equation
  yields coupling polynomial; M\"{o}bius map on the state
  parameter; normalisation by
  range~\cite{Nelson2026polynomial} \\
2 (this note) & AQUAL field equation, three physical
  postulates, single-channel detailed balance \\
3 (companion note) & Quadratic tetrad equation, two
  parity-related roots, three boundary conditions on $\mu$
  directly~\cite{Nelson2026mobius} \\
\hline
\end{tabular*}
\vspace{6pt}

Route~1 reads the algebraic content of the tetrad equation.
Route~3 reads its solution structure. Route~2 (this note) uses
neither: it derives the constitutive law from the field equation's
equilibrium response under capacity saturation. The mathematics is
two-state kinetics, not algebra or geometry. No coupling polynomial
appears. No M\"{o}bius map is invoked. No boundary conditions on
$\mu$ are assumed; instead, $\mu$ is identified with the equilibrium
occupancy, and its properties follow from the balance equation.

% ----------------------------------------------------------------------

\section*{Summary}

The AQUAL field equation requires an interpolating function
$\mu(x)$ that has been undetermined since 1984. Within the minimal
single-channel saturating closure class:

\begin{enumerate}
\item Postulate~2 (deep-field scaling) forces the driving exponent
  $n = 1$.
\item Detailed balance with linear driving and linear saturation
  yields the equilibrium $\epsilon = x/(1+x)$.
\item Scale absorption removes the rate ratio, leaving no free
  parameters.
\item Postulate~1 (Newtonian recovery) is satisfied automatically.
\end{enumerate}

The result is
\begin{equation}
  \mu(x) = \frac{x}{x + 1},
\end{equation}
the unique equilibrium response of a single-channel,
capacity-limited gravitational field under detailed balance.
What Milgrom~\cite{Milgrom1983} discovered empirically and Famaey
and Binney~\cite{FamaeyBinney2005} selected by fitting is the
unique solution within the minimal closure class consistent with
the observed asymptotic limits.

% ----------------------------------------------------------------------

\begin{thebibliography}{99}

\bibitem{Nelson2026geometry}
S.~Nelson,
\textit{The Coordinate-Space Geometry of the Tetrad Frame},
Class.\ Quantum Grav.\ (2026), submitted.

\bibitem{Nelson2026polynomial}
S.~Nelson,
\textit{Supplementary Material: The Constitutive Law from the
Coupling Polynomial} (2026).

\bibitem{Nelson2026mobius}
S.~Nelson,
\textit{Supplementary Material: The Constitutive Law from the
Quadratic Tetrad Equation} (2026).

\bibitem{BekensteinMilgrom1984}
J.~D.~Bekenstein and M.~Milgrom,
\textit{Does the missing mass problem signal the breakdown of
Newtonian gravity?}
Astrophys.\ J.\ \textbf{286} (1984) 7.

\bibitem{Milgrom1983}
M.~Milgrom,
\textit{A modification of the Newtonian dynamics as a possible
alternative to the hidden mass hypothesis},
Astrophys.\ J.\ \textbf{270} (1983) 365.

\bibitem{FamaeyBinney2005}
B.~Famaey and J.~Binney,
\textit{Modified Newtonian dynamics in the Milky Way},
Mon.\ Not.\ R.\ Astron.\ Soc.\ \textbf{363} (2005) 603.

\end{thebibliography}

\end{document}
